\chapter{面向弱监督分割与配准的医学影像持续学习研究}
\label{chap:task-level-continual-learning}
\setlength{\emergencystretch}{1em}

\section{引言}
\label{sec:task-cl-introduction}

前述章节从持续医学图像分割的场景与评价出发，建立了任务序列、历史访问和能力维度之间的对应关系，并进一步研究了多中心环境下无回放联邦持续医学图像分类。上述研究表明，任务持续演化并不是唯一需要处理的变化。模型在每个阶段能够获得的监督形式，以及后续训练能否重新调用历史样本，同样决定了学习问题的性质。尤其在医学图像分割和配准中，密集标注或成对训练数据的获取成本较高，当前任务监督受限与历史任务持续到达往往会同时出现。

本章关注两类输出形式不同、但均受到训练信息访问条件约束的医学影像任务。第一类是 ScribbleCL。该研究将第三章的 Domain-CL、Class-CL 和 Organ-CL 扩展到当前阶段仅提供涂鸦、后续阶段不回放历史原始图像的条件。当前任务监督采用 Zhang 等提出的 ZScribbleSeg\cite{zhang2026zscribbleseg}。在 Domain-CL 和 Organ-CL 中，本文将其与 EWC 结合并记为 ZS-Regu；在 Class-CL 中，将其与 MiB 结合并记为 ZS-MiB。ZScribbleSeg、EWC 和 MiB 的已有机制均按原始文献引用，不作为本论文新增算法。ScribbleCL 的新增范围是三类持续分割设定下的数据集构建、方法组合、实验设计与实证分析。

第二类是持续医学图像配准。与分割预测离散语义标签不同，配准模型需要从固定图像和移动图像估计连续空间变换。当脑部 MRI、腹部 CT、肺部 CT 和腹部 MR--CT 等任务依次到达时，图像外观、解剖区域、模态关系和损失定义均可能发生变化。SAMCL 采用有限经验回放保存少量历史图像对，并通过元持续学习协调当前样本与历史样本的更新方向；在此基础上，锐度感知优化进一步约束参数邻域中的最坏损失，以改善顺序训练中的跨任务泛化\cite{wang2024samcl,foret2021sam}。因此，SAMCL 属于有限回放的集中式持续学习方法，不能表述为无回放方法，也不提供形式化隐私保证。

本章并不试图以同一算法统一分割与配准。二者共享的研究主线是：在任务序列持续演化时，如何根据当前监督和历史访问条件构造可验证的学习目标。涂鸦分割同时面对当前阶段监督覆盖不足与历史知识保持，持续配准则在有限历史回放条件下关注旧任务保持、新任务适应和未见任务泛化。基于这一边界，本章依次给出问题定义、ZScribbleSeg 弱监督基础、三类 ScribbleCL 方法与实验设定，以及 SAMCL 方法与实验，最后比较两类任务的监督形式、历史访问、输出空间和评价差异。

\section{问题定义与研究动机}
\label{sec:task-cl-formulation}

\subsection{涂鸦监督分割及其持续扩展}
\label{subsec:task-cl-scribble-formulation}

沿用第\ref{chap:foundations}章的分割记号，设医学图像为
$\mathbf{I}:\Omega\rightarrow\mathbb{R}^{C}$，分割模型
$f_{\theta}$ 在位置 $\mathbf{u}\in\Omega$ 输出类别概率
$\mathbf{p}_{\theta}(\mathbf{u})$。在涂鸦监督条件下，只有位置集合
$\Omega_L\subset\Omega$ 具有类别标签，未标注集合为
$\Omega_U=\Omega\setminus\Omega_L$。式\eqref{eq:foundations-partial-ce}
给出的部分交叉熵只在 $\Omega_L$ 上产生直接监督，避免把未知像素简单视为背景。然而，当 $|\Omega_L|\ll|\Omega|$ 时，已标注位置不能直接给出目标的完整覆盖、边界和连通关系。模型即使在涂鸦位置上预测正确，也可能在未标注区域形成欠分割、边界偏移或碎片化结果。

静态涂鸦监督学习在固定训练分布上优化单一模型。设训练集为
$\mathcal{D}^{\mathrm{scr}}=\{(\mathbf{I}_i,\mathbf{Y}^{\mathrm{scr}}_i)\}_{i=1}^{N}$，其中
$\mathbf{Y}^{\mathrm{scr}}_i$ 只在 $\Omega_{L,i}$ 上定义。一般目标可以写为
\begin{equation}
\widehat{\theta}
=
\operatorname*{arg\,min}_{\theta}
\frac{1}{N}\sum_{i=1}^{N}
\left[
\mathcal{L}_{\mathrm{PCE}}^{(i)}(\theta)
+
\mathcal{R}_{\mathrm{scr}}^{(i)}(\theta)
\right],
\label{eq:scribble-static-objective}
\end{equation}
其中，$\mathcal{R}_{\mathrm{scr}}$ 表示从输入变换、预测一致性或空间结构中获得的附加约束。该目标的关键不是把未标注像素补写为确定标签，而是在保持监督边界的同时，为未标注区域引入可解释且受条件限制的训练信号。

ScribbleCL 进一步考虑任务按阶段到达。设任务序列为
$\mathcal{T}=\{1,\ldots,T\}$，第 $t$ 个阶段获得涂鸦训练集
$\mathcal{D}^{\mathrm{scr}}_t$，新增类别集合为 $\mathcal{C}_t$，截至该阶段的已见类别集合为
$\mathcal{C}_{1:t}=\bigcup_{\tau=1}^{t}\mathcal{C}_{\tau}$。完成第 $t$ 个阶段后，模型参数记为 $\theta_t$。令
$d_{t,j}$ 表示模型 $\theta_t$ 在任务 $j$ 测试集上的平均 Dice，则
$\mathbf{D}^{\mathrm{scr}}=[d_{t,j}]$ 仍可按照第\ref{chap:benchmark}章的阶段--任务性能矩阵进行组织。与完整监督 Class-CL 不同，此时主对角线 $d_{t,t}$ 同时受到持续约束和涂鸦监督不足的影响，旧任务条目则同时反映历史保持和早期弱监督模型本身的上限。因此，只报告最终平均 Dice 无法判断性能损失主要来自当前任务没有学好，还是来自后续训练遗忘了早期任务。

ScribbleCL 在后续阶段不重新读取历史原始图像。Domain-CL 和 Organ-CL 允许保留上一阶段参数及 EWC 所需的参数重要性统计，Class-CL 允许保留上一阶段模型以完成 MiB 蒸馏。当前阶段的弱监督目标由 ZScribbleSeg 提供，持续学习部分则根据三类场景的输出变化分别采用参数正则或背景语义建模。完整的阶段--任务性能矩阵用于区分当前任务监督不足与后续顺序更新造成的历史性能下降。

\subsection{持续医学图像配准}
\label{subsec:task-cl-registration-formulation}

医学图像配准的基本映射与训练目标已在式\eqref{eq:foundations-reg-learning}中给出。设第 $t$ 个配准任务的数据集为
$\mathcal{D}^{(t)}=\{(I_{F,i}^{(t)},I_{M,i}^{(t)})\}_{i=1}^{N_t}$，配准网络
$g_{\theta}$ 输出从固定图像域到移动图像域的变换
$\phi_{\theta}^{(t)}=g_{\theta}(I_F^{(t)},I_M^{(t)})$。不同任务可以具有不同模态、解剖区域和监督信息，因此其配准损失记为
$\mathcal{L}^{(t)}(\phi_{\theta};I_F,I_M)$。若所有已见任务数据能够同时访问，则完成前 $K$ 个任务后的理想累积目标为
\begin{equation}
\widehat{\theta}^{(K)}
=
\operatorname*{arg\,min}_{\theta}
\sum_{t=1}^{K}
\mathbb{E}_{(I_F,I_M)\sim\mathcal{D}^{(t)}}
\left[
\mathcal{L}^{(t)}(\phi_{\theta};I_F,I_M)
\right],
\qquad K=1,\ldots,T.
\label{eq:samcl-cumulative-objective}
\end{equation}
顺序训练时，当前阶段通常只能完整访问 $\mathcal{D}^{(K)}$。若直接在新任务上继续优化，参数可能偏向当前分布，并降低旧任务上的解剖重叠或地标对应精度。

SAMCL 采用容量受限的记忆缓冲区 $\mathcal{M}$ 保存少量历史图像对，满足
$|\mathcal{M}|\leq b$。这一访问条件介于联合训练与无回放学习之间：模型不能重新调用全部历史数据，但后续更新仍会读取缓冲区中的历史原始图像对。有限回放用于提供旧任务梯度，元学习用于减少当前梯度与历史梯度之间的干扰，锐度感知优化则关注参数邻域而非单个参数点的损失。最终模型既要保持已见任务的配准能力，也希望对尚未参与训练的任务形成更好的直接泛化，但这些目标只能在给定任务顺序、记忆预算和评价数据上进行验证。

表\ref{tab:task-cl-problem-comparison}概括了本章两类研究的主要差别。该比较用于明确问题边界，而不是把 Dice、目标配准误差和形变质量压缩为同一数值。

\begin{table}[htbp]
\centering
\caption{本章弱监督持续分割与持续配准问题的主要差异}
\label{tab:task-cl-problem-comparison}
\small
\renewcommand{\arraystretch}{1.18}
\setlength{\tabcolsep}{4.0pt}
\resizebox{\textwidth}{!}{%
\begin{tabular}{p{2.2cm}p{3.2cm}p{3.0cm}p{3.0cm}p{3.6cm}}
\toprule
研究对象 & 输出空间 & 当前阶段监督 & 历史访问 & 主要评价问题 \\
\midrule
ScribbleCL &
像素或体素级离散语义标签 &
少量涂鸦位置具有直接类别标签，并使用被引用的弱监督机制扩展训练信号 &
不回放历史原始图像；Domain-CL 和 Organ-CL 保留参数重要性统计，Class-CL 保留上一阶段模型 &
当前任务建模、历史知识保持、遗忘、前向能力及最终全类别分割 \\
\addlinespace
SAMCL &
固定图像域到移动图像域的连续空间变换 &
按任务使用图像相似性、解剖标签或地标构造损失 &
容量受限的历史图像对经验回放 &
已见任务配准精度、后向能力、形变合理性与未见任务泛化 \\
\bottomrule
\end{tabular}}
\end{table}

\section{ScribbleCL：面向涂鸦监督的持续医学图像分割研究}
\label{sec:scribblecl}

ScribbleCL 研究稀疏监督与任务顺序变化同时存在时的医学图像分割。当前阶段仅提供由密集标签模拟得到的涂鸦，后续训练不重新读取历史原始图像。该研究沿用第三章的 Domain-CL、Class-CL 和 Organ-CL 数据组织与评价方式，并把 ZScribbleSeg 作为当前任务弱监督基础\cite{zhang2026zscribbleseg}。Domain-CL 和 Organ-CL 使用 ZS-Regu，Class-CL 使用 ZS-MiB。两种方法分别组合已有的 EWC 与 MiB，不改变这些方法的原始学术归属\cite{kirkpatrick2017ewc,cermelli2020mib}。

\subsection{ZScribbleSeg 弱监督分割基础}
\label{subsec:prior-scribble-foundation}

设医学图像为 $\mathbf{I}:\Omega\rightarrow\mathbb{R}^{C}$，涂鸦标签为
$\mathbf{Y}^{\mathrm{scr}}$。只有位置集合 $\Omega_L\subset\Omega$ 具有可用标签，其他位置构成未标注集合 $\Omega_U=\Omega\setminus\Omega_L$。部分交叉熵只在已标注位置计算：
\begin{equation}
\mathcal{L}_{\mathrm{PCE}}
=
-\frac{1}{|\Omega_L|}
\sum_{\mathbf{u}\in\Omega_L}
\log p_{\theta}
\left(
y_{\mathbf{u}}\mid\mathbf{I},\mathbf{u}
\right).
\label{eq:scribblecl-pce}
\end{equation}
该损失避免把未标注像素直接视为背景，但监督梯度只来自少量位置。涂鸦覆盖不足、类别标注比例与真实区域比例不一致，以及缺少边界和连通结构信息，均可能限制分割结果。

\subsubsection{监督增强与全局一致性}
\label{subsubsec:prior-scribble-supervision-augmentation}

ZScribbleSeg 通过图像--涂鸦同步混合和随机遮挡扩大单次更新中的有效监督范围\cite{zhang2026zscribbleseg}。混合操作同时作用于输入图像和涂鸦标签，避免图像内容与稀疏标签错位。随机遮挡进一步改变前景与背景涂鸦的位置分布。为限制增强过程造成的预测偏移，方法比较“先预测再变换”和“先变换再预测”的结果，并以余弦距离构造全局一致性损失 $\mathcal{L}_{\mathrm{global}}$。该项约束模型对既定增强的响应一致性，但不把未标注像素直接转换为确定标签。

\subsubsection{类别比例、空间先验与结构约束}
\label{subsubsec:prior-scribble-spatial-shape}

涂鸦中各类别的像素频率通常不能代表完整目标的面积比例。ZScribbleSeg 在已标注与未标注像素的类条件外观分布近似一致的条件下，使用期望最大化估计未标注区域中的类别混合比例\cite{zhang2026zscribbleseg}。估计结果用于控制各类别候选区域的覆盖范围。

为确定具体像素的类别归属，方法结合空间距离与图像强度构造局部相似性，并据此对未标注像素进行类内排序。类别比例确定候选覆盖范围，排序后剩余的像素被用于构造较可信的负集合，形成空间先验损失 $\mathcal{L}_{\mathrm{spatial}}$。对满足单一连通结构假设的目标，方法还利用最大连通区域构造结构损失 $\mathcal{L}_{\mathrm{shape}}$。对于可能包含多个离散病灶的类别，该结构项不应直接启用。

\subsubsection{总体目标与原始论文结果}
\label{subsubsec:prior-scribble-training-objective}

ZScribbleSeg 的总体训练目标为
\begin{equation}
\begin{aligned}
\mathcal{L}_{\mathrm{ZS}}
={}&
\mathcal{L}_{\mathrm{PCE}}
+
\lambda_1\mathcal{L}_{\mathrm{global}}
\\
&+
\lambda_2\mathcal{L}_{\mathrm{spatial}}
+
\lambda_3\mathcal{L}_{\mathrm{shape}}.
\end{aligned}
\label{eq:prior-scribble-total-loss}
\end{equation}
其中，部分交叉熵提供直接监督，全局一致性用于约束增强前后的预测关系，空间先验利用类别比例与局部关系筛选未标注区域，结构损失用于减少不合理的碎片化预测。原始实现先使用部分交叉熵、全局一致性和结构损失进行预热，再启用空间先验。空间排序和类别比例估计只在训练阶段执行，不改变测试阶段的网络前向过程\cite{zhang2026zscribbleseg}。

Zhang 等比较了骨架涂鸦、随机游走、定向随机游走和随机点。原文结果显示，在控制标注像素数时，更分散的涂鸦对不规则病灶更有帮助。例如，在 MyoPS 上，骨架涂鸦的平均 Dice 为 0.281，定向随机游走和随机点分别为 0.588 和 0.595。随机点需要更多独立绘制操作，因此该结果主要说明监督覆盖与空间随机性的重要性，而不表示随机点在实际标注中始终具有更低成本\cite{zhang2026zscribbleseg}。

原论文在 ACDC 验证集上的消融显示，仅使用部分交叉熵时，平均 Dice 和 HD 分别为 0.813 和 69.58 mm；完整 ZScribbleSeg 的对应结果为 0.899 和 8.70 mm。监督增强、全局一致性、类别比例、空间先验和结构约束表现出互补作用，但这些数值来自固定数据划分下的静态弱监督实验，不能直接说明持续任务中的历史知识保持\cite{zhang2026zscribbleseg}。

原论文进一步在 ACDC、MSCMRseg、BTCV、MyoPS，以及 Medical Segmentation Decathlon（MSD）的 BrainTumor 和 Prostate 两个任务上进行评价。相同骨干下，MSCMRseg 的平均 Dice 从 PCE 的 0.385 提高到 0.870，MyoPS 从 0.281 提高到 0.636。该方法在不同结构和维度的数据上取得了具有竞争力的结果，但并未在所有数据集和所有指标上形成统一最优值。因此，这些实验仅用于说明本章当前任务弱监督模块的来源和适用范围\cite{zhang2026zscribbleseg}。

\subsection{面向三类持续分割设定的方法}
\label{subsec:scribblecl-methods}

设任务序列为 $\mathcal{T}=\{1,\ldots,T\}$，第 $t$ 个阶段获得涂鸦训练集
$\mathcal{D}^{\mathrm{scr}}_t$，完成该阶段训练后的模型参数为 $\theta_t$。所有方法只读取当前阶段原始图像，不保存或回放历史原始图像。三类场景使用不同的输出组织，因此持续学习组件也分别设计。

\subsubsection{顺序训练参照}
\label{subsubsec:scribblecl-sequential-baselines}

Dense-Sequential 使用当前阶段的密集标签直接顺序更新模型，不加入专门的历史知识保护。该方法与第三章的 Non-CL 完全对应，本节仅通过重命名强调其监督形式。PCE-Sequential 使用式\eqref{eq:scribblecl-pce}训练当前任务。ZS-Sequential 使用式\eqref{eq:prior-scribble-total-loss}训练当前任务。后两者使用完全相同的涂鸦数据，区别仅在于是否使用 ZScribbleSeg 的监督增强、类别比例和结构约束。

\subsubsection{Domain-CL 与 Organ-CL：ZS-Regu}
\label{subsubsec:scribblecl-zs-regu}

Domain-CL 保持前列腺标签语义和共享输出头不变，变化主要来自中心分布。Organ-CL 的器官、模态和输出目标依次改变，但各任务通过独立输出头组织。两类场景均不存在 Class-CL 中由新增类别引起的背景语义变化，因此采用 EWC 约束共享参数\cite{kirkpatrick2017ewc}。

完成任务 $t-1$ 后，令 $F_i^{t-1}$ 为参数 $\theta_i$ 的 Fisher 重要性估计，在线累计重要性写为
\begin{equation}
\overline{F}_i^{t-1}
=
\gamma\overline{F}_i^{t-2}
+
F_i^{t-1},
\qquad
\overline{F}_i^{0}=0.
\label{eq:scribblecl-online-fisher}
\end{equation}
ZS-Regu 在第 $t$ 个阶段优化
\begin{equation}
\mathcal{L}^{(t)}_{\mathrm{ZS\mbox{-}Regu}}
=
\mathcal{L}^{(t)}_{\mathrm{ZS}}
+
\frac{\lambda_{\mathrm{EWC}}}{2}
\sum_{i\in\Theta_{\mathrm{sh}}}
\overline{F}_i^{t-1}
\left(
\theta_i-\theta^{\star}_{t-1,i}
\right)^2,
\label{eq:scribblecl-zs-regu}
\end{equation}
其中，$\Theta_{\mathrm{sh}}$ 表示跨阶段共享的参数，$\theta^{\star}_{t-1}$ 为上一阶段训练结束时的参数。Domain-CL 对共享主干和共享分割头施加该约束。Organ-CL 对共享主干施加约束，非当前任务的输出头不参与当前阶段损失。该设计只保留参数和重要性统计，不重新读取历史图像。

\subsubsection{Class-CL：ZS-MiB}
\label{subsubsec:scribblecl-zs-mib}

Class-CL 逐步扩展统一输出空间。当前阶段只标注新增类别时，旧类别和未来类别可能被包含在当前背景中。ZS-MiB 在 ZScribbleSeg 当前任务损失中引入 MiB 的背景语义处理和旧模型蒸馏\cite{cermelli2020mib}。

令 $\mathcal{C}_{1:t-1}$ 为旧类别集合，$\mathcal{C}_t$ 为当前新增类别集合，$b$ 表示背景。对当前阶段的涂鸦背景标签，定义
\begin{equation}
\overline{p}_{t}(y\mid\mathbf{u})
=
\begin{cases}
p_t(y\mid\mathbf{u}),
& y\in\mathcal{C}_t,\\[2mm]
\displaystyle
p_t(b\mid\mathbf{u})
+
\sum_{c\in\mathcal{C}_{1:t-1}}
p_t(c\mid\mathbf{u}),
& y=b.
\end{cases}
\label{eq:scribblecl-mib-supervision-probability}
\end{equation}
背景感知的部分交叉熵为
\begin{equation}
\mathcal{L}^{(t)}_{\mathrm{uPCE}}
=
-\frac{1}{|\Omega_L^{(t)}|}
\sum_{\mathbf{u}\in\Omega_L^{(t)}}
\log
\overline{p}_{t}
\left(
y_{\mathbf{u}}^{(t)}\mid\mathbf{u}
\right).
\label{eq:scribblecl-unbiased-pce}
\end{equation}
该项在当前标签为背景时汇总旧类别与背景概率，避免把当前未标注的旧类别全部压向单一背景输出。

蒸馏时，旧模型的各旧类别概率分别与当前模型对应类别匹配。旧模型的背景概率则与当前模型的背景及新增类别概率之和匹配：
\begin{equation}
\widetilde{p}_{t}(c\mid\mathbf{u})
=
\begin{cases}
p_t(c\mid\mathbf{u}),
& c\in\mathcal{C}_{1:t-1},\\[2mm]
\displaystyle
p_t(b\mid\mathbf{u})
+
\sum_{k\in\mathcal{C}_{t}}
p_t(k\mid\mathbf{u}),
& c=b.
\end{cases}
\label{eq:scribblecl-mib-distillation-probability}
\end{equation}
相应蒸馏损失为
\begin{equation}
\mathcal{L}^{(t)}_{\mathrm{MiB}}
=
-\frac{1}{|\Omega|}
\sum_{\mathbf{u}\in\Omega}
\sum_{c\in\mathcal{C}_{1:t-1}\cup\{b\}}
p_{t-1}(c\mid\mathbf{u})
\log
\widetilde{p}_{t}(c\mid\mathbf{u}).
\label{eq:scribblecl-mib-distillation}
\end{equation}
ZS-MiB 将式\eqref{eq:prior-scribble-total-loss}中的
$\mathcal{L}_{\mathrm{PCE}}$ 替换为
$\mathcal{L}^{(t)}_{\mathrm{uPCE}}$，其余 ZScribbleSeg 项只使用当前阶段可访问的涂鸦和预测。最终目标为
\begin{equation}
\mathcal{L}^{(t)}_{\mathrm{ZS\mbox{-}MiB}}
=
\mathcal{L}^{(t)}_{\mathrm{ZS,u}}
+
\lambda_{\mathrm{MiB}}
\mathcal{L}^{(t)}_{\mathrm{MiB}},
\label{eq:scribblecl-zs-mib}
\end{equation}
其中，$\mathcal{L}^{(t)}_{\mathrm{ZS,u}}$ 表示使用背景感知部分交叉熵后的 ZScribbleSeg 当前任务损失。该方法只保留上一阶段模型，不使用历史原始图像。

\begin{table}[htbp]
\centering
\caption{ScribbleCL 比较方法及其适用场景}
\label{tab:scribblecl-methods}
\small
\renewcommand{\arraystretch}{1.15}
\setlength{\tabcolsep}{5.5pt}
\resizebox{\textwidth}{!}{%
\begin{tabular}{lccc}
\toprule
方法 & 当前阶段监督 & 历史知识处理 & 适用场景 \\
\midrule
Dense-Sequential & 密集标签 & 无 & Domain-CL、Class-CL、Organ-CL \\
PCE-Sequential & 涂鸦与部分交叉熵 & 无 & Domain-CL、Class-CL、Organ-CL \\
ZS-Sequential & 涂鸦与 ZScribbleSeg & 无 & Domain-CL、Class-CL、Organ-CL \\
ZS-Regu & 涂鸦与 ZScribbleSeg & EWC 参数正则 & Domain-CL、Organ-CL \\
ZS-MiB & 涂鸦与 ZScribbleSeg & MiB 背景建模与蒸馏 & Class-CL \\
\bottomrule
\end{tabular}}
\end{table}

\subsection{实验设定}
\label{subsec:scribblecl-experimental-settings}

\subsubsection{数据集、任务序列与数据划分}
\label{subsubsec:scribblecl-datasets}

ScribbleCL 完全沿用第\ref{chap:benchmark}章的三个数据场景、任务顺序和数据划分\cite{wang2026benchmark}。每个任务均随机划分为 60\% 训练集、15\% 验证集和 25\% 测试集。涂鸦只由训练集密集标签生成，验证集和测试集继续使用密集标签。构建涂鸦后不重新划分样本，确保 Dense-Sequential 与各弱监督方法使用相同的训练、验证和测试样本。

\begin{table}[htbp]
\centering
\caption{ScribbleCL 三类持续分割场景的数据组织}
\label{tab:scribblecl-datasets}
\scriptsize
\renewcommand{\arraystretch}{1.16}
\setlength{\tabcolsep}{3.0pt}
\resizebox{\textwidth}{!}{%
\begin{tabular}{p{1.7cm}p{6.6cm}p{3.2cm}p{2.4cm}}
\toprule
场景 & 数据集与任务顺序 & 输出组织 & 训练/验证/测试 \\
\midrule
Domain-CL &
NCI-ISBI 2013 的 A、B 中心，I2CVB 的 C 中心，以及 PROMISE12 的 D、E、F 中心，按 A$\rightarrow$B$\rightarrow$C$\rightarrow$D$\rightarrow$E$\rightarrow$F 到达 &
共享二通道输出头 &
60\%/15\%/25\% \\
\addlinespace
Class-CL &
MMWHS 心脏 CT，依次学习 LV、LA、MYO，随后加入 RV、RA，最后加入 AA、PA &
共享八通道输出头 &
60\%/15\%/25\% \\
\addlinespace
Organ-CL &
LAScarQS 左心房 LGE MRI $\rightarrow$ PROMISE12 中心 D 前列腺 MRI $\rightarrow$ LiTS 肝脏 CT $\rightarrow$ FeTS 2021 脑肿瘤 FLAIR MRI &
每个任务独立二通道输出头 &
60\%/15\%/25\% \\
\bottomrule
\end{tabular}}
\end{table}

Domain-CL 的 C 中心执行裁剪和空间对齐，所有轴位切片调整为
$256\times256$，并使用非零体素的均值和标准差进行归一化。Class-CL 将 MMWHS 图像的平面分辨率重采样为
$0.78\times0.78$ mm，平均层厚为 1.60 mm，随后调整为
$256\times256$。Organ-CL 对每幅图像进行零均值、单位方差归一化，并将轴位切片调整为 $256\times256$。以上处理均与第三章保持一致。

\subsubsection{涂鸦数据集构建}
\label{subsubsec:scribblecl-scribble-construction}

训练集涂鸦采用 WSL4MIS 开源实现中的
\texttt{scribbles\_generator.py} 构建\cite{luo2022scribble}\footnote{\url{https://github.com/HiLab-git/WSL4MIS/blob/main/code/scribbles_generator.py}}。对每个类别，先从阶段可访问的密集标签中得到二值掩膜，再逐轴位切片生成骨架。面积较大的区域先进行随机腐蚀，随后使用 Lee 骨架化获得中心线。前景骨架经过连通区域筛选和随机分支裁剪，并进行一次形态学膨胀。对面积超过 1000 像素的目标，脚本在原类别区域内施加 $[-6,6]$ 像素的随机平移和不超过 $15^\circ$ 的随机旋转，以增加涂鸦形态的变化。

生成结果保留前景类别和背景涂鸦，其他位置写为 255，并在部分交叉熵中作为忽略位置。若训练代码内部使用其他 \texttt{ignore\_index}，只能在数据读取阶段执行确定性映射，不改变磁盘中构建的数据含义。所有比较方法共享同一份固定涂鸦文件，避免由重新随机生成标注造成额外差异。

Class-CL 先按照当前阶段的标签空间完成类别重映射，再调用涂鸦生成代码。当前不可访问的旧类别和未来类别只按该阶段的背景含义处理，生成过程不读取其真实类别身份。完整密集标签仅用于构建模拟涂鸦、验证、测试和 Dense-Sequential，不进入 PCE-Sequential、ZS-Sequential 或 ZS-MiB 的训练损失。

\subsubsection{训练与评价设置}
\label{subsubsec:scribblecl-training-evaluation}

所有实验采用 ResUNet32。每个任务训练 150 个 epoch，优化器为 SGD，初始学习率为 0.008，在第 80 个 epoch 后乘以 0.5，batch size 为 8。Dense-Sequential 使用第三章对应的密集监督顺序训练设置。PCE-Sequential、ZS-Sequential、ZS-Regu 和 ZS-MiB 使用相同网络、任务顺序、优化器和训练预算。ZScribbleSeg 的损失权重设置为
$\lambda_1=0.05$、$\lambda_2=1$ 和 $\lambda_3=1$。结构项只对满足单一连通结构假设的类别启用；对于多灶或天然多连通目标，通过组件门控关闭该项。在本章统一的 150-epoch 训练预算下，前 15 个 epoch 不启用空间先验，自第 16 个 epoch 起加入空间先验损失。上述设置在所有采用 ZScribbleSeg 的比较方法之间保持一致\cite{zhang2026zscribbleseg}。

ZS-Regu 使用 $\lambda_{\mathrm{EWC}}=1$ 和 $\gamma=0.1$。ZS-MiB 使用蒸馏权重 10 和 $\alpha=1.0$。这些持续学习参数与第三章保持一致。所有方法均不回放历史原始图像，也不增加随任务保存的原始样本缓存。

评价使用第三章定义的阶段--任务性能矩阵。Domain-CL 报告 A-Dice、BWTR、RMA 和 E-FWT；Class-CL 报告 A-Dice、BWTR、RMA 和 WCD；Organ-CL 报告 A-Dice、BWTR 和 RMA。所有指标均在密集测试标签上计算。Dense-Sequential 的数值直接采用第三章 Non-CL 结果，其余方法在实验完成前不填写数值。

\subsection{实验结果}
\label{subsec:scribblecl-results}

\begin{table}[htbp]
\centering
\caption{Domain-CL 场景下 ScribbleCL 方法的结果}
\label{tab:scribblecl-domain-results}
\small
\renewcommand{\arraystretch}{1.14}
\setlength{\tabcolsep}{7pt}
\begin{tabular}{lcccc}
\toprule
方法 & A-Dice & BWTR & RMA & E-FWT \\
\midrule
Dense-Sequential & $.676\pm.021$ & $-.237\pm.038$ & $1.086\pm.008$ & $.260\pm.151$ \\
PCE-Sequential & -- & -- & -- & -- \\
ZS-Sequential & -- & -- & -- & -- \\
ZS-Regu & -- & -- & -- & -- \\
\bottomrule
\end{tabular}
\vspace{2pt}

\begin{minipage}{0.96\textwidth}
\footnotesize 注：Dense-Sequential 数值直接采用第\ref{chap:benchmark}章 Domain-CL 的 Non-CL 结果；“--”表示相应实验结果尚未填入。
\end{minipage}
\end{table}

\begin{table}[htbp]
\centering
\caption{Class-CL 场景下 ScribbleCL 方法的结果}
\label{tab:scribblecl-class-results}
\small
\renewcommand{\arraystretch}{1.14}
\setlength{\tabcolsep}{7pt}
\begin{tabular}{lcccc}
\toprule
方法 & A-Dice & BWTR & RMA & WCD \\
\midrule
Dense-Sequential & $.471\pm.012$ & $-.598\pm.015$ & $1.014\pm.028$ & $.572\pm.023$ \\
PCE-Sequential & -- & -- & -- & -- \\
ZS-Sequential & -- & -- & -- & -- \\
ZS-MiB & -- & -- & -- & -- \\
\bottomrule
\end{tabular}
\vspace{2pt}

\begin{minipage}{0.96\textwidth}
\footnotesize 注：Dense-Sequential 数值直接采用第\ref{chap:benchmark}章 Class-CL 的 Non-CL 结果；“--”表示相应实验结果尚未填入。
\end{minipage}
\end{table}

\begin{table}[htbp]
\centering
\caption{Organ-CL 场景下 ScribbleCL 方法的结果}
\label{tab:scribblecl-organ-results}
\small
\renewcommand{\arraystretch}{1.14}
\setlength{\tabcolsep}{10pt}
\begin{tabular}{lccc}
\toprule
方法 & A-Dice & BWTR & RMA \\
\midrule
Dense-Sequential & $.603\pm.008$ & $-.369\pm.005$ & $1.006\pm.008$ \\
PCE-Sequential & -- & -- & -- \\
ZS-Sequential & -- & -- & -- \\
ZS-Regu & -- & -- & -- \\
\bottomrule
\end{tabular}
\vspace{2pt}

\begin{minipage}{0.88\textwidth}
\footnotesize 注：Dense-Sequential 数值直接采用第\ref{chap:benchmark}章 Organ-CL 的 Non-CL 结果；“--”表示相应实验结果尚未填入。
\end{minipage}
\end{table}

\subsection{实验分析}
\label{subsec:scribblecl-experimental-analysis}

\subsection{讨论}
\label{subsec:scribblecl-discussion}

\subsection{结论}
\label{subsec:scribblecl-conclusion}
\section{SAMCL：锐度感知元持续医学图像配准方法}
\label{sec:samcl}

SAMCL 面向解剖区域、成像模态和配准关系依次变化的顺序任务。该方法以有限经验回放为历史信息载体，以元持续学习协调当前与历史梯度，并用锐度感知目标改善参数邻域中的损失。其研究目标是获得一个能够在给定任务序列上保持多任务配准能力、并对后续未见任务具有一定直接泛化能力的单模型\cite{wang2024samcl}。

\subsection{面向配准任务序列的元持续学习}
\label{subsec:samcl-meta-continual}

直接顺序训练只根据当前任务梯度更新参数，容易使模型离开旧任务的低损失区域。经验回放从容量为 $b$ 的缓冲区
$\mathcal{M}$ 中采样历史图像对，并与当前样本交替更新。仅把两类损失相加仍可能产生梯度冲突。Meta Experience Replay 的核心思想是通过一阶元学习，使来自不同样本或任务的梯度方向更加一致。对从缓冲区采样的批次
$z_i=(I_{F,i},I_{M,i})$，其近似目标可以写为
\begin{equation}
\min_{\theta}
\mathbb{E}_{z_1,\ldots,z_b\sim\mathcal{M}}
\left[
\sum_{i=1}^{b}
\left(
\mathcal{L}^{(t_i)}_i(\theta)
-
\frac{\alpha}{2}
\sum_{j=1}^{i-1}
\nabla_{\theta}\mathcal{L}^{(t_i)}_i(\theta)^{\top}
\nabla_{\theta}\mathcal{L}^{(t_j)}_j(\theta)
\right)
\right].
\label{eq:samcl-meta-objective}
\end{equation}
当两个梯度内积为正时，沿一个样本下降也更可能降低另一任务的损失；内积为负则提示干扰。式\eqref{eq:samcl-meta-objective}通过元更新鼓励跨样本梯度对齐，从而在有限记忆下近似协调当前和历史任务。该机制依赖缓冲区对历史分布的代表性，不能等同于访问完整历史数据。

SAMCL 使用水库采样维护记忆，使顺序到达的图像对具有进入固定容量缓冲区的机会。每个当前样本先触发当前任务内层更新，随后从记忆中采样历史图像对执行回放更新，再进行一阶元更新。该流程使当前任务和历史任务在同一参数轨迹上产生作用。记忆中的样本是原始固定--移动图像对，因此方法的有效历史访问条件是有限回放，而非仅保存参数统计。

\subsection{锐度感知优化与当前任务--历史任务联合更新}
\label{subsec:samcl-sharpness-update}

可变形配准目标通常具有高度非凸性。即使某一参数点在当前批次上损失较低，其邻域内损失仍可能快速上升，导致模型对采样噪声或新分布敏感。SAMCL 将元持续学习中的任务损失替换为锐度感知损失
\begin{equation}
\mathcal{L}^{(t)}_{\mathrm{sam}}(\theta)
=
\max_{\|\boldsymbol{\epsilon}\|_2\leq\rho}
\mathcal{L}^{(t)}(\theta+\boldsymbol{\epsilon}),
\label{eq:samcl-sharpness-loss}
\end{equation}
其中，$\rho$ 控制参数邻域半径。外层优化不只降低参数点 $\theta$ 的损失，还降低该邻域中的近似最坏损失，从而偏向相对平坦的局部区域\cite{foret2021sam}。

对内层最大化在 $\boldsymbol{\epsilon}=0$ 处作一阶展开，可以得到扰动近似
\begin{equation}
\widehat{\boldsymbol{\epsilon}}(\theta)
=
\rho
\frac{
\nabla_{\theta}\mathcal{L}^{(t)}(\theta)
}{
\|\nabla_{\theta}\mathcal{L}^{(t)}(\theta)\|_2
}.
\label{eq:samcl-perturbation}
\end{equation}
实际更新使用扰动参数
$\theta+\widehat{\boldsymbol{\epsilon}}(\theta)$ 处的梯度。把
$\mathcal{L}^{(t)}_{\mathrm{sam}}$ 代入式\eqref{eq:samcl-meta-objective}后，当前样本和回放样本均采用锐度感知更新，元更新再协调两类轨迹。该组合针对的是损失几何和跨样本梯度关系：经验回放提供历史证据，元学习减小任务干扰，锐度感知优化改善局部邻域的稳健性。它并不直接约束形变的生物力学真实性，形变平滑与可逆性仍由配准模型和正则项决定。

SAMCL 采用 LK-UNet 作为配准骨干，首层使用 16 个卷积核。配准变换由平稳速度场参数化，并采用对称微分同胚积分；膜能与弯曲能共同约束形变平滑性。每个任务训练 $10000$ 次迭代，优化器为 Adam，学习率为 $10^{-4}$，批大小为 $4$。记忆容量为 $200$，元学习率为 $0.25$，锐度邻域半径为 $0.05$。这些设置用于四任务顺序实验，其作用范围受任务顺序、骨干和记忆预算限制。


\IfFileExists{figures/ch05/samcl_framework.pdf}{%
\begin{figure}[htbp]
\centering
\includegraphics[width=\textwidth]{ch05/samcl_framework.pdf}
\caption{SAMCL 面向顺序配准任务的训练与部署框架。训练阶段依次接收脑部 MR--MR、腹部 CT--CT、肺部 CT--CT 和腹部 MR--CT 图像对，并通过有限经验回放、元更新和锐度感知优化形成单一配准模型。}
\label{fig:samcl-framework}
\end{figure}
}{}

\subsection{配准持续学习实验}
\label{subsec:samcl-experiments}

\subsubsection{跨器官与跨模态持续配准任务序列}
\label{subsubsec:samcl-protocol}

实验按照 OASIS 脑部 MR--MR、腹部 CT--CT、NLST 肺部 CT--CT 和腹部 MR--CT 的顺序训练。四个任务同时包含解剖区域、模态关系和配准类型变化。OASIS 与腹部 CT--CT 为个体间配准，NLST 和腹部 MR--CT 为个体内配准。表\ref{tab:samcl-datasets}给出数据规模、重采样和评价指标。所有体数据裁剪或填充为
$112\times96\times112$，OASIS 重采样到 $2\times2\times2$ mm，其余任务重采样到 $3\times3\times3$ mm。CT 强度裁剪到 $[-200,300]$ 后，所有输入线性归一化到 $[0,1]$；实验不使用额外数据增强。

\begin{table}[htbp]
\centering
\caption{SAMCL 四个顺序配准任务的数据组织与评价}
\label{tab:samcl-datasets}
\small
\renewcommand{\arraystretch}{1.14}
\setlength{\tabcolsep}{4.2pt}
\resizebox{\textwidth}{!}{%
\begin{tabular}{llllll}
\toprule
顺序 & 数据集与任务 & 配准关系 & 训练/验证/测试 & 重采样间距 & 评价 \\
\midrule
1 & OASIS，脑部 MR--MR & 个体间 & 290/40/84 & $2\times2\times2$ mm & 35 类平均 Dice \\
2 & Abdomen CT--CT & 个体间 & 21/3/6 & $3\times3\times3$ mm & 13 类平均 Dice \\
3 & NLST，肺部 CT--CT & 个体内 & 294/42/84 & $3\times3\times3$ mm & 地标 TRE(mm) \\
4 & Abdomen MR--CT & 个体内、跨模态 & 80/9/16 & $3\times3\times3$ mm & 9 类平均 Dice \\
\bottomrule
\end{tabular}}
\end{table}

OASIS、腹部 CT--CT 和腹部 MR--CT 的任务损失由负局部归一化互相关与软 Dice 损失构成，局部相关窗口大小为 $3$。NLST 使用负局部归一化互相关与基于对应地标的距离损失。由于四个任务可用的解剖标签和地标不同，任务损失并不完全相同；本章只在各任务自身的评价尺度上比较顺序训练前后的变化，不把不同任务的原始 Dice 和 TRE 直接求平均。

比较方法包括未配准输入、独立任务训练、集中式多任务训练、直接顺序训练，以及 EWC、SI、GPM、MER 和 SAMCL。独立训练为每个任务使用单独模型，集中式多任务训练同时访问全部任务数据，因此二者与持续方法的历史访问条件不同，只作为任务特定和集中访问参考。EWC、SI 和 GPM 不回放原始历史图像；MER 与 SAMCL 使用有限回放。

令 $R_{i,j}$ 表示完成任务 $i$ 后在任务 $j$ 上的评价。实验对每个任务报告最终性能
$\mathrm{AVG}_j=R_{T,j}$ 和后向变化
$\mathrm{BWT}_j=R_{T,j}-R_{j,j}$。对 Dice，数值越高越好，负 BWT 表示性能下降；对 TRE，数值越低越好，正 BWT 表示误差增大。因此，表中 Dice-BWT 应越接近零或越高越好，TRE-BWT 则越低越好。最后一个任务之后没有后续训练，其 BWT 按定义为零。

\subsubsection{配准精度、知识保持、遗忘与泛化分析}
\label{subsubsec:samcl-analysis}

表\ref{tab:samcl-main-results}转录四任务顺序结束后的主要结果。直接顺序训练在前三个任务上出现较明显的知识损失：OASIS 和腹部 CT--CT 的 Dice-BWT 分别为 $-0.395$ 和 $-0.129$，NLST 的 TRE-BWT 为 $3.508$ mm。SAMCL 将三项后向变化改善为 $-0.037$、$-0.039$ 和 $0.306$ mm，说明有限回放与元持续更新能够缓解直接顺序训练的遗忘。

\begin{table}[htbp]
\centering
\caption{SAMCL 四任务顺序实验的最终性能与后向变化。OASIS、Abd CT 和 Abd MR--CT 报告 Dice（越高越好），NLST 报告 TRE（越低越好）。}
\label{tab:samcl-main-results}
\scriptsize
\renewcommand{\arraystretch}{1.12}
\setlength{\tabcolsep}{3.0pt}
\resizebox{\textwidth}{!}{%
\begin{tabular}{lcccccccc}
\toprule
& \multicolumn{4}{c}{最终性能 $R_{T,j}$} & \multicolumn{4}{c}{后向变化 $R_{T,j}-R_{j,j}$} \\
\cmidrule(lr){2-5}\cmidrule(lr){6-9}
方法 & OASIS & Abd CT & NLST TRE & Abd MR--CT & OASIS & Abd CT & NLST TRE & Abd MR--CT \\
\midrule
未配准 & 0.598 & 0.311 & 7.932 & 0.353 & -- & -- & -- & -- \\
独立训练 & 0.738 & 0.335 & 3.259 & 0.550 & -- & -- & -- & -- \\
集中式多任务 & 0.743 & 0.358 & 2.662 & 0.506 & -- & -- & -- & -- \\
直接顺序训练 & 0.348 & 0.270 & 6.922 & 0.568 & $-0.395$ & $-0.129$ & 3.508 & 0 \\
EWC & 0.659 & 0.364 & 3.836 & 0.506 & $-0.083$ & $-0.011$ & 0.386 & 0 \\
SI & 0.461 & 0.303 & 8.186 & 0.555 & $-0.281$ & $-0.084$ & 4.685 & 0 \\
GPM & 0.458 & 0.290 & 7.095 & 0.556 & $-0.281$ & $-0.085$ & 3.574 & 0 \\
MER & 0.698 & 0.331 & 3.611 & 0.532 & $-0.036$ & $-0.040$ & 0.077 & 0 \\
SAMCL & 0.697 & 0.352 & 3.855 & 0.563 & $-0.037$ & $-0.039$ & 0.306 & 0 \\
\bottomrule
\end{tabular}}
\end{table}

与不使用回放的 SI 和 GPM 相比，SAMCL 在 OASIS 与腹部 CT--CT 上取得更高的最终 Dice，并具有更小的后向下降。与 EWC 相比，SAMCL 在 OASIS 上具有更高的最终 Dice 和更小的后向下降；在腹部 CT--CT 上，SAMCL 的最终 Dice 略低于 EWC，且 EWC 的后向下降更小。与同样使用有限回放的 MER 相比，SAMCL 的表现更接近：OASIS 的最终 Dice 和前两个任务的 BWT 基本相当，SAMCL 在腹部 CT--CT 与最后到达的 MR--CT 上取得更高 Dice，而 MER 在 NLST 上具有更低的最终 TRE 和更小的 TRE-BWT。该结果不支持“在所有任务上统一优于 MER”的表述，更合适的解释是锐度感知元持续学习在多项任务上保持了竞争性，并改变了不同任务之间的稳定性--可塑性取舍。

SAMCL 在腹部 MR--CT 上取得 0.563 的最终 Dice，高于独立训练的 0.550 和集中式多任务训练的 0.506，但直接顺序训练在最后任务上为 0.568。最后任务没有后续遗忘，因此该结果主要反映当前任务适应，而不能用于说明历史保持。OASIS 和腹部 CT--CT 上，SAMCL 仍低于集中式多任务或独立训练参考，说明有限记忆不能完全恢复完整历史访问条件。

记忆容量消融显示，在 OASIS、腹部 CT--CT 和腹部 MR--CT 上，较大的缓冲区总体上提高了最终配准精度，反映历史样本覆盖对有限回放的重要性。锐度感知消融则比较了带有和不带有 SAM 的元持续模型。对于腹部 CT--CT 到腹部 MR--CT 的同域或相关任务，以及腹部 CT--CT 到 NLST、OASIS 到 NLST 等跨解剖或跨域任务，带有 SAM 的模型在原始图中均获得更好的未见任务点估计。该观察支持锐度感知训练对给定未见任务泛化的作用，但实验只覆盖少量预定义迁移方向，不能外推为任意模态或器官上的通用保证。

\IfFileExists{figures/ch05/samcl_qualitative.pdf}{%
\begin{figure}[htbp]
\centering
\includegraphics[width=\textwidth]{ch05/samcl_qualitative.pdf}
\caption{不同训练策略在 OASIS 与腹部 CT 图像对上的代表性配准结果。}
\label{fig:samcl-qualitative}
\end{figure}
}{}

\IfFileExists{figures/ch05/samcl_ablation_generalization.pdf}{%
\begin{figure}[htbp]
\centering
\includegraphics[width=\textwidth]{ch05/samcl_ablation_generalization.pdf}
\caption{SAMCL 的记忆容量消融与未见任务泛化分析。左图中的 Dice 与 TRE 方向不同，右图覆盖预设的有限迁移方向。}
\label{fig:samcl-ablation-generalization}
\end{figure}
}{}

实验采用固定任务顺序和单一主要骨干，未系统交换任务到达次序，也未覆盖更广泛的成像设备、解剖区域或配准损失。有限回放需要保存历史原始图像对，其存储规模和治理要求不同于无回放方法。表中给出测试样本上的点估计与配对检验标记，但未报告多随机种子训练的均值和标准差。因此，本节结论限定于该四任务序列与既定实现。

\section{综合讨论}
\label{sec:task-cl-discussion}

本章两项研究对应不同的训练信息访问条件。ScribbleCL 在当前阶段只获得涂鸦，并在后续训练中不回放历史原始图像。Domain-CL 和 Organ-CL 通过 ZS-Regu 保存参数重要性统计，Class-CL 通过 ZS-MiB 保存上一阶段模型。SAMCL 的当前任务图像对相对完整，但历史任务通过容量受限的原始图像对重新参与训练。两类研究需要分别报告当前监督、历史信息载体和可访问范围。

ScribbleCL 所采用的 ZScribbleSeg 机制作用于当前任务的稀疏监督，EWC 和 MiB 分别处理参数保持与类别增量中的背景语义变化。三者均为已有方法\cite{zhang2026zscribbleseg,kirkpatrick2017ewc,cermelli2020mib}。本论文在三类 Benchmark 场景中完成数据集构建和组合评价，不把静态弱监督结果写成持续学习结果。SAMCL 则通过有限回放、元持续更新和锐度感知优化处理顺序配准，两部分不共享同一算法。

SAMCL 则在同一训练循环中同时使用当前和历史图像对。经验回放提供历史梯度，元更新减少跨样本干扰，锐度感知目标偏向邻域损失较平坦的参数。该组合可以共同影响稳定性与可塑性，但仍受记忆容量限制。更大的缓冲区可能改善旧任务覆盖，同时增加存储和治理成本；较强的历史约束也可能降低新任务适应。原始四任务结果中，SAMCL 在若干 Dice 任务上优于 MER，但在 NLST 的 TRE 和后向变化上不及 MER，正体现了不同任务之间的取舍，而不是单一指标上的统一优势。


\section{本章小结}
\label{sec:task-cl-summary}

本章研究了任务持续演化与训练信息受限在医学图像分割和配准中的两种具体表现。对于涂鸦监督持续分割，本章构建 ScribbleCL，并沿用第三章的 Domain-CL、Class-CL 和 Organ-CL 数据场景。当前任务监督采用已有的 ZScribbleSeg\cite{zhang2026zscribbleseg}；Domain-CL 和 Organ-CL 使用 ZS-Regu，Class-CL 使用 ZS-MiB。实验同时设置 Dense-Sequential、PCE-Sequential 和 ZS-Sequential，用于区分监督形式与持续知识处理的作用。上述方法组合不改变 ZScribbleSeg、EWC 和 MiB 的原始学术归属。

对于持续医学图像配准，本章介绍了 SAMCL。该方法在容量受限的经验回放基础上使用一阶元持续学习协调当前与历史梯度，并将锐度感知损失应用于两类更新，以改善参数邻域中的稳健性。四个顺序配准任务表明，SAMCL 相较直接顺序训练明显缓解了早期任务性能下降，并在多个任务上与 MER、独立训练和集中式多任务参考保持竞争性；同时，NLST 结果、记忆依赖和单一任务顺序也限定了其结论范围。通过上述两类研究，本章进一步说明，医学影像持续学习方法必须同时报告任务如何变化、当前监督提供什么信息、历史数据以何种形式可访问，以及相应结果能够支持何种范围的结论。
